% !TeX root = Report.tex
% Introduction body. The chapter heading is declared in Report.tex.
\enlargethispage{2\baselineskip}

Задача обнаружения прямых линий на изображении и оценивания их параметров
возникает во многих задачах обработки изображений и компьютерного зрения. 
Обнаружение дорожной разметки на изображениях с автомобильных камер, 
выделение границ объектов, анализ медицинских или спутниковых снимков --- это лишь
небольшая часть приложений, в которых требуется надежно выделять прямые.

Одним из основных решений этой задачи является
преобразование Хафа \cite{duda1972hough}. Его идея состоит в переходе от
множества точек изображения к пространству параметров прямых, на которых
лежат эти точки: при параметризации прямой через пару $(\rho, \theta)$
каждая точка изображения порождает кривую в пространстве параметров. Если
несколько точек принадлежат одной и той же прямой, то соответствующие им
кривые пересекаются в общей точке. На этом принципе строится аккумуляторная
функция, зависящая от $(\rho, \theta)$; ее значения отражают число точек
изображения, лежащих на прямой с заданными параметрами. Поэтому оценка
параметров искомой прямой сводится к нахождению точки максимума аккумуляторной
функции.

Однако преобразование Хафа чувствительно к
шуму на входном изображении: шумовые точки порождают лишние кривые в
пространстве параметров, которые могут влиять на максимум аккумуляторной
функции и, следовательно, на точность оценивания соответствующих параметров.
Поэтому применение данного преобразования требует предварительной обработки
изображения.

В данной работе в качестве специальной предобработки рассматриваются алгоритмы
на основе метода анализа сингулярного спектра SSA и его комплекснозначного
варианта Complex SSA (CSSA) \cite{ts_structure,ssa_with_r,zhornikova2016}.
Сначала к строкам изображения применяется дискретное преобразование Фурье,
затем к полученным строкам или к столбцам комплексной матрицы
применяется CSSA, после чего выполняется обратное преобразование Фурье по
строкам или столбцам. Применение подхода на основе SSA после преобразования Фурье
рассматривалось ранее \cite{trickett2003}. В данной работе исследуются два
основных варианта такой предобработки: \texttt{CSSA ROW.ROW}, в котором CSSA и обратное DFT применяются по строкам, и
\texttt{CSSA COL.ROW}, в котором  CSSA применяется по столбцам, а обратное DFT по строкам. Кроме того, рассматриваются
два способа улучшения предобработки на основе CSSA: ESPRIT
\cite{ssa_with_r} и пороговая фильтрация восстановленного изображения.

Для сравнения с подходами на основе CSSA используются более простые и стандартные методы
предобработки: медианный фильтр \cite{huang1979}, фильтр Винера
\cite{wiener1949,lim1990}, а также метод \texttt{MAX.ROW}, в котором в каждой
строке сохраняются несколько наибольших элементов. Сравнительный обзор первых двух фильтров
для подавления разных видов шума на изображениях приведен в
\cite{mafi2019}, где основное внимание уделяется качеству
восстановленного изображения в целом. В настоящей работе, в силу особенностей
преобразования Хафа, важно также, сколько посторонних шумовых элементов
остается после предобработки и как они влияют на оценивание параметров прямых.
Это делает актуальным сравнение методов в качестве предобработки перед применением преобразования Хафа.

Цель работы состоит в исследовании методов предобработки изображений
для последующего обнаружения прямых преобразованием Хафа и в сравнении этих
методов по точности оценивания параметров прямых.

Для достижения этой цели в работе решаются следующие задачи:
\begin{enumerate}
	\setlength{\itemsep}{0pt}
	\setlength{\parsep}{0pt}
	\setlength{\topsep}{0pt}
	\item описать модель изображения, преобразование Хафа, метод SSA, его
	комплекснозначный вариант CSSA и используемые стандартные методы
	предобработки;
	\item исследовать варианты \texttt{CSSA ROW.ROW} и \texttt{CSSA COL.ROW},
	включая точное восстановление и свойства ранга
	траекторных матриц;
	\item проанализировать особенности восстановления после предобработки на основе CSSA,
	их влияние на преобразование Хафа, а также способы их уменьшения с помощью
	ESPRIT и пороговой фильтрации;
	\item провести численное сравнение методов предобработки по точности оценок
	$\rho$ и $\theta$ после применения преобразования Хафа в случаях одной и
	нескольких прямых.
\end{enumerate}

Структура работы такова. В первой главе описываются используемые алгоритмы,
вводится модель изображения и фиксируются обозначения. Во второй главе
исследуется использование CSSA для предобработки изображений: отдельно
рассматриваются варианты \texttt{ROW.ROW} и \texttt{COL.ROW}, обсуждаются
ранговые свойства соответствующих рядов, а также рассматриваются ESPRIT и
пороговая фильтрация. В третьей главе приводится численное сравнение методов
предобработки по точности оценивания параметров прямых преобразованием Хафа, в
том числе анализируется влияние дискретизации пространства параметров,
уровня шума и числа прямых. В заключении подводятся основные итоги работы и
намечаются возможные направления дальнейшего исследования.
